\documentclass[11pt,letterpaper,dvipsnames]{article}

\usepackage{style} 
%\usepackage{soul}
\addbibresource{refs.bib}

\title{An Exponential Top Fan-In Hierarchy for Border Depth-\texorpdfstring{$4$}{4} Circuits}
\author{
    \href{https://anakin.phd/}{Anakin Dey} \thanks{Department of Mathematics, The Ohio State University~\href{mailto:dey.92@osu.edu}
    {{\ttfamily\upshape dey.92@osu.edu}}; partially supported by NSF award DMS-2231565.} \and 
}

\newif\ifmynotes
% \mynotestrue
\mynotesfalse
\newcommand{\ADnote}[1]{\ifmynotes\textcolor{orange}{[Anakin: #1]}\else\fi}

\begin{document}

\date{}
\maketitle

\begin{abstract}
  Wow how abstract and not concrete.
\end{abstract}

%\pranjal{Check the abstract and title!} \ADnote{I think both look fine.}

\newpage 
\tableofcontents

\section{Introduction}\label{sec:intro}

\begin{thrm}[Impossibility, see~{\Cref{thrm:impossibility}}]\label{thrm:intro-exact}
    For every fixed $k, \delta \geq 1$, if $n > 2(k + \delta)\delta^{2k}$, then $P_{k + 1, n}$ cannot be computed by any $\Sigma^{[k]}\Pi\Sigma\Pi^{[\delta]}$ circuit, regardless of its size.
\end{thrm}
We prove the theorem above in \Cref{sec:impossibility}.
Thus, depth-$4$ circuits with top fan-in $k + 1$ and bottom fan-in $1$ are already more powerful than depth-$4$ circuits with top fan-in $k$ and any constant bottom fan-in. 
This is because $P_{k + 1, n}$ can be computed by $\Sigma^{[k+1]}\Pi\Sigma$ circuits of size $\bigO{kn}$, but no $\Sigma^{[k]}\Pi\Sigma\Pi^{[\delta]}$ circuit computes $P_{k + 1, n}$. 
For fixed $\delta$, this threshold is at most singly exponential in $k$.
We remark that even showing the result for homogeneous circuits is nontrivial, and we prove it as a first step in \Cref{thrm:impossibility-exact}.


\subsection{Commutative Algebra}\label{sec:prelim-ca}

We state basic definitions and results that can also be found in standard commutative algebra texts~\cite{Eis95, AM69, Mat86, Lan02} and computational algebra texts~\cite{vzGG13, CLO25}.

For $S \subseteq R$, write $\ideal{S}$ for the ideal generated by $S$. 
The \emph{radical} $\sqrt{I}$ of an ideal $I$ is 
\[
    \set{a \in R | a^n \in I \text{ for some } n \geq 1}.
\]
An ideal $I$ is \emph{radical} if $I = \sqrt{I}$.
A proper ideal $\mathfrak{p}$ is \emph{prime} if $R/\mathfrak{p}$ is an integral domain and \emph{maximal} if $R/\mathfrak{p}$ is a field.
Maximal ideals are prime, and prime ideals are radical.
A \emph{minimal prime} of $I$ is a prime ideal minimal among those containing $I$.

A ring is \emph{Noetherian} if every ideal is finitely generated.
If $R$ is Noetherian, then so are its quotients, its localizations, finitely generated $R$-algebras, and formal power-series rings over $R$ in finitely many variables~\cite[Propositions~7.1 and~7.3, Corollary~7.7, and Corollary~10.27]{AM69}.
In a Noetherian ring, every prime ideal containing an ideal $I$ contains a minimal prime of $I$.
Moreover, every proper ideal in a Noetherian ring has finitely many minimal primes, whose intersection is its radical~\cite[Proposition~4.6 and Theorem~7.13]{AM69}.

\subsection{Algebraic Geometry}\label{sec:prelim-ag}

We give basic definitions and results for the algebro-geometric tools we use.
Much of this can be found in standard textbooks~\cite{Har77, Sha13, CLO25, BCS97, HP94} or online sources such as the Stacks project~\cite{stacks-project}.

Throughout \Cref{sec:prelim-ag}, let $\F$ be an algebraically closed field unless otherwise stated.
The $N$-dimensional \emph{affine space} $\A^N$ over $\F$ is $\F^N$.
For $S \subseteq \F[\vec{x}]$, let $\V(S) \subseteq \A^N$ be its common zero set; for $T\subseteq\A^N$, let $\I(T)$ be the ideal of polynomials vanishing on $T$. 
A \emph{hypersurface} is the zero set of a nonconstant polynomial.
The sets $\V(S)$ are the \emph{affine varieties} and form the closed sets of the \emph{Zariski topology}.
We allow varieties to be reducible or empty.
The \emph{Zariski closure} $\overline{T}$ of $T$ is $\V(\I(T))$, the smallest closed subset of $\A^N$ containing $T$.
A nonempty variety is \emph{irreducible} if it is not the union of two proper closed subsets. 
Every affine variety $V$ is a finite union $V_1 \cup \cdots \cup V_r$ of irreducible affine varieties with $V_i \not\subseteq V_j$ for $i \neq j$.
These varieties are the \emph{irreducible components} of $V$, and the decomposition is unique up to order~\cite[Chapter~4, \S6, Theorem~4]{CLO25}.

An \emph{affine open subset} of a variety is an open subset isomorphic to an affine variety.
A property holds \emph{generally} or \emph{generically} on a variety $V$ if it holds on a dense open subset of $V$.
In an irreducible variety, every nonempty open subset is dense and irreducible, and every finite intersection of nonempty open subsets is nonempty~\cite[Chapter~I, Exercise~1.6]{Har77}.
In particular, if $V$ is irreducible and $g_1,\ldots,g_r\notin\I(V)$, then $\set{b \in V | g_1(b) \cdots g_r(b) \neq 0}$ is a nonempty dense open subset of $V$.

For $S \subseteq [N]$, write $\A^S$ for the affine space with coordinates $\vec{x}_S$ and $\pi_S \colon \A^N \to \A^S$ for the coordinate projection $a \mapsto (a_s)_{s \in S}$.
For $T \subseteq [N]$ and $S \subseteq \F[\vec{x}_T]$, the notation $\V(S)$ denotes the zero set of $S$ in $\A^N$, not $\A^T$, unless otherwise noted.

For every ideal $J \subseteq \F[\vec{x}]$, the \emph{Nullstellensatz} gives 
\[ 
    \I(\V(J)) = \sqrt{J}, \qquad \V(J) = \emptyset \iff J = \F[\vec{x}]. 
\]
Thus $\V$ and $\I$ give mutually inverse, inclusion-reversing bijections between radical ideals of $\F[\vec{x}]$ and affine varieties in $\A^N$, with prime ideals corresponding to irreducible varieties~\cite[Chapter~4, \S2, Theorems~6--7; \S5, Corollary~4]{CLO25}.
For every proper ideal $J$, the irreducible components of $\V(J)$ are precisely $\V(\mathfrak p)$ as $\mathfrak p$ ranges over the minimal primes of $J$~\cite[Corollary~I.1.4 and Proposition~I.1.5]{Har77}.

For an affine variety $V$, define its \emph{coordinate ring} by $\F[V] \defeq \F[\vec{x}]/\I(V)$.
For $f \in \F[\vec{x}]$, define its restriction to $V$ by $f|_V \defeq f + \I(V)$.
If $V$ is irreducible, then its \emph{function field} $\F(V)$ is $\mathrm{Frac}(\F[V])$.
The coordinates $\vec{x}_S$ are \emph{algebraically independent on $V$} if the family $(x_s|_V)_{s \in S}$ is algebraically independent over $\F$.
The coordinates $\vec{x}_S$ are algebraically independent on $V$ precisely when $\I(V) \cap \F[\vec{x}_S] = \set{0}$.
In this case we identify $\F[\vec{x}_S]$ with its image in $\F[V]$.

\begin{lem}\label{lem:units}
    Let $V \subseteq \A^N$ be an affine variety and let $q \in \F[\vec{x}]$.
    The restriction $q|_V$ is a unit of $\F[V]$ if and only if $V \cap \V(q) = \emptyset$.
\end{lem}
\begin{pf}
    If $q|_V$ is a unit, then evaluating an inverse at each point of $V$ shows that $q$ does not vanish there.
    Conversely, if $\V(\I(V) + \ideal{q}) = \emptyset$, then the Nullstellensatz gives $1 = a + qh$ for some $a \in \I(V)$ and $h \in \F[\vec{x}]$.
    Restricting to $V$ gives $(q|_V)(h|_V) = 1$.
\end{pf}

\section{Conclusion}\label{sec:conclusion}

\paragraph{Acknowledgments}

\paragraph{AI Disclosure}


% STOC says put bibliography at very end
\printbibliography

\end{document}

